%==============================================================================
% MUWS 2026 draft #3 — Why is multimodal hate harder than sarcasm?
%==============================================================================
\documentclass[sigconf,nonacm]{acmart}
\settopmatter{printacmref=false}
\usepackage{booktabs}
\usepackage{graphicx}

\begin{document}

\title{Why Is Multimodal Hate Harder than Sarcasm?\\
A Matched-Budget Adaptation Study with Small Vision--Language Models}

\author{Anonymous Author(s)}
\affiliation{\institution{Anonymous Institution}\country{}}

\begin{abstract}
Hateful-meme and multimodal-sarcasm detection are routinely benchmarked side by side, but
their relative difficulty for compact models is rarely controlled. We hold the model
(SmolVLM-500M), adaptation method (LoRA), and budget (3{,}000 examples, 2 epochs) fixed
and compare the two tasks. Under an identical recipe, sarcasm improves by $+0.28$ AUROC
($0.600\to0.881$) while hate improves by only $+0.04$ ($0.581\to0.619$)---a $7\times$
difference in adaptation gain. We analyse the gap: training loss converges to $\sim$0.1
for both tasks, yet held-out hate barely moves, indicating that the model fits
hate-training quirks that do not transfer to the hate dev set---the signature of the
benign-confounder design of Hateful Memes. We argue that ``multimodal toxicity'' should
not be treated as a single difficulty class, with implications for how small-model
moderation is benchmarked and reported.
\end{abstract}

\keywords{hateful memes, multimodal sarcasm, task difficulty, LoRA, small VLMs}

\maketitle

\section{Introduction}
Papers on multimodal human understanding frequently evaluate hate and sarcasm detection
with a shared pipeline and report them as comparable tasks. We ask a controlled question:
under an \emph{identical} small-model adaptation budget, how much does each task actually
improve? The answer is sharply different, and the difference is informative about what
makes multimodal hate hard.

\section{Setup}
We fine-tune SmolVLM-500M-Instruct with LoRA~\cite{hu2022lora} ($r{=}16$, 1.85\% of
parameters) on each task separately using the same generative yes/no SFT recipe, the same
3{,}000 training examples per task, 2 epochs, bf16, and identical prompts. We evaluate on
balanced splits ($n{=}500$) of Hateful Memes~\cite{kiela2020hateful} and
MMSD2.0~\cite{qin2023mmsd2}, reporting AUROC and the prediction-positive rate. We also log
training loss to relate fitting to generalisation.

\section{Results}
\paragraph{Adaptation gain differs $7\times$.} Table~\ref{tab:gain} shows that the same
recipe yields a $+0.281$ AUROC gain on sarcasm but only $+0.038$ on hate.

\begin{table}[t]
\centering
\caption{Matched-budget LoRA on SmolVLM-500M (fp16, $n{=}500$). Identical recipe; very
different gains.}
\label{tab:gain}
\begin{tabular}{lccc}
\toprule
Task & zero-shot AUROC & +LoRA AUROC & $\Delta$ \\
\midrule
Hateful Memes & 0.581 & 0.619 & $+0.038$ \\
MMSD2.0       & 0.600 & 0.881 & $+0.281$ \\
\bottomrule
\end{tabular}
\end{table}

\paragraph{Fitting is not the bottleneck.} Both runs drive training loss to $\sim$0.1,
i.e.\ the model \emph{can} memorise the supervised yes/no targets for hate just as well as
for sarcasm. Yet only sarcasm generalises to held-out data. This is consistent with the
Hateful Memes design~\cite{kiela2020hateful}: benign confounders mean that surface
patterns learnable from the train set do not characterise the dev set, so a small model
with light adaptation overfits to non-transferable cues.

\paragraph{De-biasing happens for both; separability only for sarcasm.} On both tasks the
prediction-positive rate moves from $\sim$0.95 toward the base rate ($0.36$ for hate,
$0.29$ for sarcasm), so the model stops answering with a constant label in both cases.
The difference is that for sarcasm this de-biasing is accompanied by genuine class
separation, whereas for hate it is not.

% TODO Fig.1: (a) train-loss curves overlaid (both ->~0.1); (b) AUROC bars with deltas;
% (c) optional qualitative confounder examples the model misclassifies.
\begin{figure}[t]
\centering
\fbox{\parbox{0.9\columnwidth}{\centering\vspace{1.2cm}\textit{Figure 1 (placeholder):
overlaid training-loss curves (both converge) beside held-out AUROC gains (only sarcasm
rises).}\vspace{1.2cm}}}
\caption{Fitting vs.\ generalisation under a matched budget.}
\label{fig:loss}
\end{figure}

\section{Discussion}
The practical message is that ``multimodal toxicity detection'' bundles tasks of very
different learnability for small models. Reporting a single small-model moderation result
across hate and sarcasm can therefore be misleading: a recipe that looks effective on
sarcasm may be near-useless on hate. For deployment, hate detection with compact models
likely needs more data, harder negatives, or explicit confounder-aware training rather
than simply more fine-tuning steps.

\section{Limitations}
Single backbone, single seed, $n{=}500$ slices, and one adaptation budget; we do not yet
sweep training-set size, so we cannot rule out that much larger hate-training data closes
the gap. We infer overfitting from the loss/held-out mismatch rather than a full
attribution study.

\section{Conclusion}
Under a matched small-model LoRA budget, multimodal sarcasm is roughly $7\times$ more
adaptable than hateful-meme detection, and the hate gap stems from generalisation, not
fitting. Multimodal toxicity should be benchmarked as a set of distinct-difficulty tasks,
not a single one.

\bibliographystyle{ACM-Reference-Format}
\bibliography{refs}
\end{document}
